% Standalone submission table; content retained from the verified manuscript.
\documentclass[fleqn,12pt]{wlpeerj}
% Shared table support. Compile Table_01.tex--Table_24.tex from the full source-package root.
% Uses the manuscript class, fonts, frozen display-value macros, and resolved citation labels.
% No scores or bibliography identities are recalculated here.
\usepackage{array,ragged2e,tabularx,longtable,makecell,adjustbox}
\usepackage[hidelinks]{hyperref}
\usepackage{float,amssymb,colortbl,multirow,fontspec,xeCJK}
\geometry{letterpaper,left=2.54cm,right=2.54cm,top=2.54cm,bottom=2.54cm,headheight=14pt}
\setmainfont{TeX Gyre Termes}
\providecommand{\TableFontPath}{fonts/noto-cjk/}
\setCJKmainfont[Path=\TableFontPath,BoldFont=CNSSSubmissionSerif-Bold.otf,ItalicFont=CNSSSubmissionSerif-Regular.otf,ItalicFeatures={FakeSlant=0.2}]{CNSSSubmissionSerif-Regular.otf}
\setCJKsansfont[Path=\TableFontPath]{CNSSSubmissionSans-Regular.otf}
\setCJKmonofont[Path=\TableFontPath]{CNSSSubmissionSans-Regular.otf}
\newcolumntype{L}{>{\RaggedRight\arraybackslash}X}
\newcolumntype{P}[1]{>{\RaggedRight\arraybackslash}p{#1}}
\newcommand{\zh}[1]{#1}
\newcommand{\en}[1]{#1}
\setlength{\emergencystretch}{2em}
\widowpenalty=10000
\clubpenalty=10000
\brokenpenalty=10000
\raggedbottom
\pagestyle{empty}
\setlength{\parindent}{0pt}
\setlength{\parskip}{0pt}
\newcommand{\SharedResultRows}{%
\SharedResultRow{GPT-5.6 Terra}{0.676}{0.658}{0.667}{0.766}{0.654}{0.689}%
\SharedResultRow{\makecell[l]{DeepSeek-V4-Pro\\(non-thinking)}}{0.651}{0.643}{0.647}{0.761}{0.643}{0.654}%
\SharedResultRow{Claude Sonnet 4.5}{0.598}{0.597}{0.598}{0.709}{0.589}{0.614}%
\SharedResultRow{\makecell[l]{DeepSeek-V4-Pro\\(thinking)}}{0.713}{0.505}{0.591}{0.671}{0.572}{0.623}%
\SharedResultRow{Kimi K2.6}{0.581}{0.563}{0.572}{0.714}{0.557}{0.599}%
\SharedResultRow{Qwen}{0.507}{0.281}{0.361}{0.470}{0.345}{0.390}%
}
\newcommand{\QwenResultRows}{%
\QwenResultRow{No adapter}{0.507}{0.281}{0.361}{0.470}{0.345}{0.390}%
\QwenResultRow{LoRA, seed 42}{0.588}{0.434}{0.500}{0.624}{0.511}{0.480}%
\QwenResultRow{LoRA, seed 43}{0.634}{0.507}{0.563}{0.686}{0.571}{0.549}%
\QwenResultRow{LoRA, seed 44}{0.633}{0.499}{0.558}{0.676}{0.567}{0.543}%
}
\newcommand{\QwenBaselineExact}{0.361}
\newcommand{\QwenLoraExact}{0.540\pm0.035}


\definecolor{guideL}{HTML}{DED9F7}
\definecolor{guideK}{HTML}{C6ECEE}
\definecolor{guideS}{HTML}{FFF0B0}
\definecolor{guideT}{HTML}{CDEACF}
\newcommand{\gspan}[2]{{\setlength{\fboxsep}{1.4pt}\colorbox{guide#1}{\strut [#2]\,\textsf{[#1]}}}}
\newcommand{\gcontext}[1]{\textcolor{black!75}{#1}}
\newcommand{\gsource}[1]{\,\textsuperscript{\textsf{[#1]}}}
\newcommand{\ggloss}[1]{\par\emph{#1}}
% Citation labels copied verbatim from the verified manuscript aux.
\AtBeginDocument{%
\newlabel{app:c}{{}{30}{Appendix C. Supplementary Evaluation}{section*.70}{}}
\bibcite{artstein2008agreement}{{1}{2008}{{Artstein and Poesio}}{{}}}
\bibcite{bender2018datastatements}{{2}{2018}{{Bender and Friedman}}{{}}}
\bibcite{bhola2020retrieving}{{3}{2020}{{Bhola et~al.}}{{}}}
\bibcite{cao2021occupational}{{4}{2021}{{Cao et~al.}}{{}}}
\bibcite{chen2021boundary}{{5}{2021}{{Chen and Kong}}{{}}}
\bibcite{chen2024jobsdf}{{6}{2024}{{Chen et~al.}}{{}}}
\bibcite{conneau2020xlmr}{{7}{2020}{{Conneau et~al.}}{{}}}
\bibcite{council2017eqf}{{8}{2017}{{Council of the European Union}}{{}}}
\bibcite{desanto2026skilllens}{{9}{2026}{{De~Santo et~al.}}{{}}}
\bibcite{devlin2019bert}{{10}{2019}{{Devlin et~al.}}{{}}}
\bibcite{ding2023annotator}{{11}{2023}{{Ding et~al.}}{{}}}
\bibcite{ESCOv1_2_2024}{{12}{2024}{{European Commission}}{{}}}
\bibcite{escopediaCompetence}{{13}{nda}{{European Commission}}{{}}}
\bibcite{escopediaSkillsPillar}{{14}{ndb}{{European Commission}}{{}}}
\bibcite{multiconer2023}{{15}{2023}{{Fetahu et~al.}}{{}}}
\bibcite{gebru2021datasheets}{{16}{2021}{{Gebru et~al.}}{{}}}
\bibcite{gnehm2022skills}{{17}{2022}{{Gnehm et~al.}}{{}}}
\bibcite{gururangan2020dapt}{{18}{2020}{{Gururangan et~al.}}{{}}}
\bibcite{hripcsak2005agreement}{{19}{2005}{{Hripcsak and Rothschild}}{{}}}
\bibcite{hu2022lora}{{20}{2022}{{Hu et~al.}}{{}}}
\bibcite{kim2026guidelines}{{21}{2026}{{Kim et~al.}}{{}}}
\bibcite{klie2024annotationquality}{{22}{2024}{{Klie et~al.}}{{}}}
\bibcite{lafferty2001crf}{{23}{2001}{{Lafferty et~al.}}{{}}}
\bibcite{levrang2014esco}{{24}{2014}{{le~Vrang et~al.}}{{}}}
\bibcite{li2026chineseskillspanpreprint}{{25}{2026}{{Li et~al.}}{{}}}
\bibcite{li2023coannotating}{{26}{2023}{{Li et~al.}}{{}}}
\bibcite{uie2022}{{27}{2022}{{Lu et~al.}}{{}}}
\bibcite{dynamicner2025}{{28}{2025}{{Luo et~al.}}{{}}}
\bibcite{macrodataPlatform}{{29}{nd}{{MacroData}}{{}}}
\bibcite{magron2024jobskape}{{30}{2024}{{Magron et~al.}}{{}}}
\bibcite{mcclelland1973competence}{{31}{1973}{{McClelland}}{{}}}
\bibcite{nguyen2024rethinking}{{32}{2024}{{Nguyen et~al.}}{{}}}
\bibcite{peng2015weibo}{{33}{2015}{{Peng and Dredze}}{{}}}
\bibcite{ILO_2020_OJVs_BigData}{{34}{2020}{{Podjanin and Strietska-Ilina}}{{}}}
\bibcite{qwen2025report}{{35}{2025}{{Qwen Team}}{{}}}
\bibcite{gollie2024}{{36}{2024}{{Sainz et~al.}}{{}}}
\bibcite{sayfullina2018softskills}{{37}{2018}{{Sayfullina et~al.}}{{}}}
\bibcite{senger2024survey}{{38}{2024}{{Senger et~al.}}{{}}}
\bibcite{tianchiRecruitment163746}{{39}{nd}{{Tianchi}}{{}}}
\bibcite{ulhaq2026nerannotation}{{40}{2026}{{Ul~Haq et~al.}}{{}}}
\bibcite{gptner2025}{{41}{2025}{{Wang et~al.}}{{}}}
\bibcite{cluener2020}{{42}{2020}{{Xu et~al.}}{{}}}
\bibcite{gliner-naacl24}{{43}{2024}{{Zaratiana et~al.}}{{}}}
\bibcite{jensen2022kompetencer}{{44}{2022a}{{Zhang et~al.}}{{}}}
\bibcite{zhang2022skillspan}{{45}{2022b}{{Zhang et~al.}}{{}}}
\bibcite{zhang2022weak}{{46}{2022c}{{Zhang et~al.}}{{}}}
\bibcite{zhang2024nnose}{{47}{2024}{{Zhang et~al.}}{{}}}
\bibcite{zhang2023escoxlmr}{{48}{2023}{{Zhang et~al.}}{{}}}
\bibcite{universalner}{{49}{2024}{{Zhou et~al.}}{{}}}
}

\begin{document}
\setcounter{table}{8}
\begingroup
\small\setlength{\tabcolsep}{4pt}\renewcommand{\arraystretch}{1.13}
% BEGIN PRESERVED TABLE
\begin{longtable}{@{}P{0.07\linewidth}P{0.49\linewidth}P{\dimexpr0.44\linewidth-16pt\relax}@{}}
\caption{Annotation principles for competency mentions and span boundaries.}\label{tab:guide-boundaries}\\
\toprule\textbf{Rule} & \textbf{Annotation principle} & \textbf{Example and English gloss}\\\midrule
\endfirsthead
\multicolumn{3}{@{}l}{\textit{Table \thetable\ (continued)}}\\
\toprule\textbf{Rule} & \textbf{Annotation principle} & \textbf{Example and English gloss}\\\midrule
\endhead


G1 & The annotation scheme distinguishes L, K, S, and T. Collapsed categories are used only in secondary evaluation.\gsource{E,P} & \gspan{L}{英语六级}\ggloss{College English Test Band 6}\par\smallskip\gspan{K}{计算机专业}\ggloss{Computer science major}\par\smallskip\gspan{S}{数据分析}\ggloss{Data analysis}\par\smallskip\gspan{T}{责任心}\ggloss{Sense of responsibility} \\[5pt]
G2 & Final annotations contain non-overlapping spans. A sentence may contain several independent mentions, with each character assigned to at most one span.\gsource{P} & \gcontext{熟悉使用}\gspan{S}{Word}\gcontext{，}\gspan{S}{Excel}\gcontext{，}\gspan{S}{PPT}\gcontext{等办公软件}\ggloss{Familiar with using Word, Excel, PowerPoint, and other office software}\par\smallskip No additional covering list span. \\[5pt]
G3 & Each mention is a contiguous source substring that preserves spelling and internal spaces. Any normalisation is stored separately.\gsource{Z,P} & \gcontext{熟悉 }\gspan{S}{Jekins}\ggloss{Familiar with Jekins}\par\smallskip Keep the source spelling, not \texttt{Jenkins}. \\[5pt]
G4 & Annotators select the shortest span that expresses a complete requirement. Words and identifiers remain intact; semantic completeness takes priority over length, with no fixed length limit.\gsource{Z,P} & \gspan{S}{支持服务}\ggloss{Support services}\par\smallskip The truncated \zh{支持服} is incomplete. \\[5pt]
G5 & Generic proficiency triggers usually remain outside the span, while heads needed to express the requirement remain inside. Excluded triggers can still inform type assignment.\gsource{Z,P} & \gcontext{会使用}\gspan{S}{Excel}\ggloss{Able to use Excel}\par\smallskip\gspan{T}{会聊天}\ggloss{Can hold a conversation}\par\smallskip\gcontext{有}\gspan{T}{服务意识}\ggloss{Has a service-oriented mindset} \\[5pt]
G6 & Coordinated items are separated when each remains meaningful in context. Shared qualifiers may stay outside the spans, and implied words are not added.\gsource{Z,P} & \gcontext{熟悉使用}\gspan{S}{Word}\gcontext{，}\gspan{S}{Excel}\gcontext{，}\gspan{S}{PPT}\gcontext{等办公软件}\ggloss{Familiar with using Word, Excel, PowerPoint, and other office software}\par\smallskip Shared-experience case: see A.3. \\*[5pt]
G7 & A necessary action, object, or shared head remains within the span when splitting would change its meaning. Span length alone does not determine type.\gsource{Z,P} & \gspan{S}{优化曝光与转化率}\ggloss{Optimise exposure and the conversion rate}\par\smallskip Isolated metric names would lose the optimisation activity. \\
\bottomrule
\end{longtable}

% END PRESERVED TABLE

% Source-key explanation is retained in the accompanying README_TABLES.md.
\endgroup
\end{document}
